% !TeX root = main_aip.tex
\newpage
\section{Manuscript Outline}


\begin{verbatim}

=========================================================
Section: Introduction
=========================================================

    Subject:
        Motivation for the development of the Field
        Particle System (FPS) and overview of its
        principal contributions.

    Purpose:
        Introduce the need for efficient and scalable
        particle dynamics frameworks, motivate the
        development of FPS, and summarize the primary
        contributions and organization of the paper.

    Paragraph Subjects:

        - Introduce particle systems as a computational
          approach for representing physical systems
          through collections of interacting particles.

        - Discuss the importance of particle interaction
          models and dynamics formulations in determining
          the behavior, accuracy, and computational
          characteristics of particle systems.

        - Identify the need for particle dynamics
          formulations that support efficient large-scale
          and massively parallel computation.

        - Introduce the Field Particle System (FPS) as a
          particle dynamics framework developed to address
          these requirements.

        - Provide a high-level overview of the FPS
          philosophy without discussing implementation
          details.

        - Summarize the principal characteristics of FPS,
          including geometric interactions,
          interpenetrating contacts, instantaneous state
          updates, energyless representation, and
          source-owned data ownership.

        - Summarize the primary contributions of the
          present work.

        - Provide an overview of the organization of the
          remainder of the manuscript.


	Paragrph 1
		Purpose: What particle dynamics methods are.
			- Defines a particle system.
			- Explains that particle motion comes from interaction rules.
			- Introduces the dominant approaches:
				* constitutive laws,
				* equations of state,
				* potential functions,
				* collision models.
			- Stops before discussing applications.
		Flow to Next:

	Paragraph 2: Why do interaction models matter


	Paragraph 3:Why look for something different?
			-  Existing methods are often phenomenon-specific.
    

	Paragraph 4:Motivation for exploring an alternative framework.

	Paragraph 5:Introduction of FPS.

	Paragraph 6:Contributions.
=========================================================
Section: Background and Related Work
=========================================================

    Subject:
        Review of particle dynamics formulations and the
        interaction models used by existing particle
        methods.

    Purpose:
        Establish the theoretical context for FPS by
        reviewing how particle interactions, momentum
        transfer, and state evolution are traditionally
        represented within particle-based systems.

    Paragraph Subjects:

        - Define the fundamental components of a particle
          dynamics system including particle state,
          interactions, and state evolution.

        - Review force-based particle dynamics and the use
          of interaction forces to govern particle motion.

        - Review energy-based and potential-based
          formulations used to determine particle
          interactions.

        - Review contact-based methods that employ overlap
          measurements and constitutive response models.

        - Review pair-owned interaction formulations and
          the update mechanisms commonly used to transfer
          momentum between particles.

        - Discuss the computational implications of pair-
          wise updates, synchronization requirements, and
          data dependencies.

        - Review particle methods developed for parallel
          computing architectures and the approaches used
          to manage interaction ownership.

        - Identify the assumptions common to existing
          particle dynamics formulations.

        - Position FPS relative to existing approaches and
          motivate investigation of an alternative
          geometric, source-owned dynamics framework.
		  
		- Review geometric interaction formulations and
		  overlap-based contact representations.

=========================================================
Section: Methodology
=========================================================

    Subject:
        Definition of the Field Particle System (FPS)
        methodology including its governing principles,
        particle representation, force generation,
        numerical integration, and source-owned
        execution model.

    Purpose:
        Present the complete FPS methodology and describe
        how geometric relationships are transformed into
        particle motion through local interactions,
        force evaluation, and numerical integration.

    Paragraph Subjects:

        - Provide a high-level overview of the FPS
          methodology and its principal computational
          stages.

        - Introduce the fundamental principles that
          govern the FPS framework.

        - Define purely geometric interactions and their
          role in determining particle behavior.

        - Define interpenetrating contacts and the
          overlap state as a finite-duration interaction
          domain.

        - Define instantaneous dynamics and the role of
          the current geometric configuration in
          determining particle response.

        - Define energyless dynamics and discuss energy
          as a diagnostic rather than a governing
          quantity.

        - Define the source-owned interaction model and
          the ownership rules governing particle updates.

        - Define geometry-to-motion causality and
          describe the conversion of geometry into
          motion through repeated local interactions.

        - Define the particle representation and the
          quantities stored by each particle.

        - Define the geometric quantities used to
          characterize particle contacts.

        - Define particle overlap and overlap area.

        - Describe construction of the contact state for
          an interacting particle pair.

        - Define the overlap-area force model.

        - Define pair stiffness and force generation.

        - Define source acceleration accumulation from
          contact contributions.

        - Describe the semi-implicit Euler integration
          procedure.

        - Define drift and excursion and discuss their
          significance in long-term system evolution.

        - Describe boundary interactions and boundary
          force generation.

        - Summarize the complete timestep execution
          sequence.

        - Discuss computational properties of the
          methodology including source ownership,
          symmetry, conservation characteristics, and
          parallel execution.

=========================================================
Section: Implementation and Performance Characteristics
=========================================================

    Subject:
        Computational implementation of FPS and its
        performance characteristics.

    Purpose:
        Describe the execution environment used for the
        demonstrations and quantify the computational
        overhead associated with the FPS dynamics model.

    Paragraph Subjects:

        - Describe the hardware and software
          configuration used for all simulations.

        - Summarize the implementation relationship
          between FPS and the underlying gPCD framework.

        - Present representative particle counts and
          simulation sizes.

        - Quantify the computational cost of FPS.

        - Compare FPS performance against the gPCD
          baseline.

        - Discuss the scalability implications of the
          observed performance characteristics.


=========================================================
Section: Simulation Demonstrations
=========================================================

    Subject:
        Representative simulations demonstrating the
        capabilities of the Field Particle System (FPS)
        across multiple physical domains.

    Purpose:
        Illustrate the range of phenomena that can be
        represented by FPS and demonstrate the emergence
        of complex behavior from the proposed dynamics
        formulation.

    Paragraph Subjects:

        - Introduce the purpose of the demonstration
          simulations and describe the visualization
          techniques used throughout this section.

        - Introduce the Cathedral simulation and its
          objectives.

        - Describe the Cathedral environment including
          the particle species used to represent air,
          dust, light, and solid boundaries.

        - Present the resulting particle behavior and
          illustrate light transport through the dust-
          filled environment.

        - Discuss the interaction of light particles
          with dust particles and reflective walls.

        - Summarize the phenomena demonstrated by the
          Cathedral simulation.

        - Introduce the Incense simulation and its
          objectives.

        - Describe the simulation configuration,
          particle species, and boundary conditions used
          to generate the rising smoke plume.

        - Present the initial laminar flow behavior
          observed near the source.

        - Present the development of periodic flow
          structures within the plume.

        - Present the transition from periodic motion
          to turbulent flow behavior.

        - Discuss the flow structures and emergent
          behavior observed throughout the transition
          process.

        - Summarize the phenomena demonstrated by the
          Incense simulation.

        - Introduce the converging-diverging nozzle
          simulation and its objectives.

        - Describe the nozzle geometry, particle
          initialization procedure, boundary conditions,
          and visualization methodology.

        - Present HSV velocity-direction visualizations
          of the resulting flow field.

        - Present acceleration of the flow through the
          converging section and throat region.

        - Present kinematic signatures consistent with
          choking behavior.

        - Present shock-like structures observed within
          the diverging section.

        - Discuss the observed flow behavior and the
          information revealed through HSV velocity-
          direction visualization.

        - Summarize the phenomena demonstrated by the
          converging-diverging nozzle simulation.

        - Summarize the collective observations from all
          demonstration cases and discuss the breadth of
          phenomena represented by FPS.
		  
		- Discuss the emergence of complex behavior from
		  a common dynamics framework.

=========================================================
Section: Conclusions
=========================================================

    Subject:
        Summary of the Field Particle System (FPS)
        methodology, its computational characteristics,
        and the phenomena demonstrated through simulation.

    Purpose:
        Summarize the principal contributions of the
        work, review the demonstrated capabilities of
        FPS, and identify directions for future
        investigation.

    Paragraph Subjects:

        - Summarize the motivation for developing FPS
          and the objectives addressed by the proposed
          methodology.

        - Summarize the principal characteristics of FPS
          including geometric interactions,
          interpenetrating contacts, source-owned
          updates, and the numerical integration
          framework.

        - Summarize the computational characteristics of
          the methodology and its suitability for
          large-scale parallel execution.

        - Summarize the Cathedral simulation and its
          demonstration of multi-species particle
          interactions including air, dust, light, and
          boundary interactions.

        - Summarize the Incense simulation and its
          demonstration of laminar, periodic, and
          turbulent flow behavior emerging from the FPS
          dynamics.

        - Summarize the converging-diverging nozzle
          simulation and the observed kinematic
          signatures consistent with choking together
          with shock-like structures.

        - Summarize the collective significance of the
          simulation demonstrations and the range of
          phenomena represented by FPS.

        - Discuss future work including expanded
          thermo-fluid modeling, additional particle
          species, larger-scale simulations, and further
          investigation of the physical phenomena
          observed in the demonstration cases.



\end{verbatim}

